%% ============================================================
%%  PLANTILLA DE APUNTES DE MATEMÁTICA
%%  Compilar con:  pdflatex plantilla_apuntes.tex   (2 veces)
%%  Para el índice analítico: makeindex plantilla_apuntes
%%  y luego pdflatex otra vez. (Con --shell-escape se hace solo.)
%% ============================================================
\documentclass[11pt,a4paper,oneside]{report}

%% ---------- Idioma y codificación ----------
\usepackage[T1]{fontenc}
\usepackage[utf8]{inputenc}
\usepackage{lmodern}
\usepackage[spanish,es-noshorthands]{babel}
\usepackage{microtype}

%% ---------- Página y párrafos ----------
\usepackage[margin=2.5cm]{geometry}
\usepackage{parskip}            % separación entre párrafos sin sangría
\usepackage{setspace}
\setstretch{1.1}
\usepackage{fancyhdr}
\pagestyle{fancy}
\fancyhf{}
\fancyhead[L]{\small\nouppercase{\leftmark}}
\fancyhead[R]{\small\thepage}
\renewcommand{\headrulewidth}{0.4pt}

%% ---------- Matemática ----------
\usepackage{amsmath,amssymb,amsthm,mathtools}
\usepackage{bm}

%% ---------- Figuras, tablas, colores ----------
\usepackage{graphicx}
\graphicspath{{figuras/}}       % pon tus imágenes en la carpeta "figuras"
\usepackage{xcolor}
\usepackage{booktabs}           % tablas bonitas (\toprule, \midrule)
\usepackage{array}
\usepackage{tabularx}
\usepackage{longtable}
\usepackage{caption}
\usepackage{float}              % permite [H] para fijar posición
\usepackage{mwe}                % solo para imágenes de ejemplo; bórralo luego
\usepackage{tikz}

%% ---------- Código ----------
\usepackage{listings}
\lstset{
  basicstyle=\ttfamily\small,
  keywordstyle=\color{blue!70!black}\bfseries,
  commentstyle=\color{green!40!black}\itshape,
  stringstyle=\color{orange!80!black},
  numbers=left, numberstyle=\tiny\color{gray}, numbersep=8pt,
  backgroundcolor=\color{gray!8},
  frame=single, rulecolor=\color{gray!40},
  breaklines=true, showstringspaces=false, tabsize=4,
  literate={á}{{\'a}}1 {é}{{\'e}}1 {í}{{\'i}}1 {ó}{{\'o}}1 {ú}{{\'u}}1
           {Á}{{\'A}}1 {É}{{\'E}}1 {Í}{{\'I}}1 {Ó}{{\'O}}1 {Ú}{{\'U}}1
           {ñ}{{\~n}}1 {Ñ}{{\~N}}1
}

%% ---------- Índices ----------
\usepackage{imakeidx}
\makeindex[title=Índice analítico, intoc]

%% ---------- Cajas de colores ----------
\usepackage[most]{tcolorbox}

%% ---------- Hipervínculos y referencias (SIEMPRE al final) ----------
\usepackage[colorlinks=true,linkcolor=blue!60!black,citecolor=green!40!black,
            urlcolor=blue!60!black]{hyperref}
\usepackage[spanish,nameinlink]{cleveref}

%% ============================================================
%%  ENTORNOS: definición, teorema, demostración, notas...
%%  Todos comparten UN solo contador, dentro de cada capítulo.
%% ============================================================
\theoremstyle{definition}
\newtheorem{definicion}{Definición}[chapter]
\newtheorem{ejemplo}[definicion]{Ejemplo}
\newtheorem{ejercicio}[definicion]{Ejercicio}

\theoremstyle{plain}
\newtheorem{teorema}[definicion]{Teorema}
\newtheorem{proposicion}[definicion]{Proposición}
\newtheorem{lema}[definicion]{Lema}
\newtheorem{corolario}[definicion]{Corolario}

\theoremstyle{remark}
\newtheorem*{observacion}{Observación}
\newtheorem*{notademo}{Nota sobre la demostración}
\newtheorem*{idea}{Idea clave}

\renewcommand{\proofname}{Demostración}
\renewcommand{\qedsymbol}{$\blacksquare$}

% Nombres para cleveref
\crefname{definicion}{definición}{definiciones}
\Crefname{definicion}{Definición}{Definiciones}
\crefname{teorema}{teorema}{teoremas}
\Crefname{teorema}{Teorema}{Teoremas}
\crefname{proposicion}{proposición}{proposiciones}
\Crefname{proposicion}{Proposición}{Proposiciones}
\crefname{lema}{lema}{lemas}
\Crefname{lema}{Lema}{Lemas}
\crefname{corolario}{corolario}{corolarios}
\Crefname{corolario}{Corolario}{Corolarios}
\crefname{ejemplo}{ejemplo}{ejemplos}
\Crefname{ejemplo}{Ejemplo}{Ejemplos}
\crefname{ejercicio}{ejercicio}{ejercicios}
\Crefname{ejercicio}{Ejercicio}{Ejercicios}
\crefname{equation}{ecuación}{ecuaciones}
\Crefname{equation}{Ecuación}{Ecuaciones}
\crefname{figure}{figura}{figuras}
\Crefname{figure}{Figura}{Figuras}
\crefname{table}{tabla}{tablas}
\Crefname{table}{Tabla}{Tablas}
\crefname{chapter}{capítulo}{capítulos}
\Crefname{chapter}{Capítulo}{Capítulos}
\crefname{section}{sección}{secciones}
\Crefname{section}{Sección}{Secciones}

% Estilo visual: barra lateral de color + fondo suave
\tcbset{
  cajabase/.style={enhanced, breakable, sharp corners, boxrule=0pt,
    frame hidden, left=8pt, right=6pt, top=4pt, bottom=4pt,
    before skip=10pt, after skip=10pt}
}
\tcolorboxenvironment{definicion}{cajabase,
  colback=blue!6,   borderline west={3pt}{0pt}{blue!65!black}}
\tcolorboxenvironment{teorema}{cajabase,
  colback=red!6,    borderline west={3pt}{0pt}{red!70!black}}
\tcolorboxenvironment{proposicion}{cajabase,
  colback=orange!7, borderline west={3pt}{0pt}{orange!80!black}}
\tcolorboxenvironment{lema}{cajabase,
  colback=orange!7, borderline west={3pt}{0pt}{orange!80!black}}
\tcolorboxenvironment{corolario}{cajabase,
  colback=orange!7, borderline west={3pt}{0pt}{orange!80!black}}
\tcolorboxenvironment{ejemplo}{cajabase,
  colback=green!5,  borderline west={3pt}{0pt}{green!45!black}}
\tcolorboxenvironment{ejercicio}{cajabase,
  colback=green!5,  borderline west={3pt}{0pt}{green!45!black}}
\tcolorboxenvironment{notademo}{cajabase,
  colback=yellow!14, borderline west={3pt}{0pt}{yellow!60!black}}
\tcolorboxenvironment{observacion}{cajabase,
  colback=gray!10,  borderline west={3pt}{0pt}{gray!70}}
\tcolorboxenvironment{idea}{cajabase,
  colback=purple!6, borderline west={3pt}{0pt}{purple!70}}

%% ============================================================
%%  ATAJOS MATEMÁTICOS
%% ============================================================
\newcommand{\N}{\mathbb{N}}
\newcommand{\Z}{\mathbb{Z}}
\newcommand{\Q}{\mathbb{Q}}
\newcommand{\R}{\mathbb{R}}
\newcommand{\C}{\mathbb{C}}
\newcommand{\eps}{\varepsilon}
\newcommand{\dd}{\mathrm{d}}
\newcommand{\dpar}[2]{\frac{\partial #1}{\partial #2}}
\DeclarePairedDelimiter{\abs}{\lvert}{\rvert}
\DeclarePairedDelimiter{\norm}{\lVert}{\rVert}
\DeclarePairedDelimiter{\pesc}{\langle}{\rangle}
\DeclarePairedDelimiter{\conj}{\{}{\}}
\DeclareMathOperator{\sen}{sen}
\DeclareMathOperator{\tg}{tg}
\DeclareMathOperator{\Img}{Im}
\DeclareMathOperator{\Ker}{Ker}
\DeclareMathOperator{\rg}{rg}

% Atajo para "término indexado": lo escribe en el texto Y lo añade al índice
\newcommand{\termino}[1]{\textbf{#1}\index{#1}}

%% ============================================================
%%  PORTADA
%% ============================================================
\title{Apuntes de Análisis}
\author{Práxedis Jiménez Ruvalcaba}
\date{\today}

\begin{document}

\maketitle
\tableofcontents

%% ============================================================
\part{Análisis I}



%% ============================================================
\part{Análisis II}

\chapter{Tercer Parcial}

\section{Examen}

\begin{ejercicio}
Sean $f,g$ funciones integrables $[a,b]$. Determine si la función 
\[ h (x) =\frac{|f(x)-g(x)|+f(x)+g(x)}{2} \]
es integrable en $[a,b]$.
\end{ejercicio}

\begin{ejercicio}
Sean $f,g$ funciones continuas en $[a,b]$ tal que para todo $x\in[a,b],f(x)>g(x)$. 
Demuestre que $\int_{a}^{b}f(x) dx > \int_{a}^{b}g(x) dx$. (Note que es estricta la desigualdad).
\end{ejercicio}

\begin{ejercicio}
Sea $f:[a,b]\rightarrow \R$ una función continua. Demuestre que existe un $c\in (a,b)$ tal que
\[ \int_{a}^{b}f(x)dx=f(c)(b-a)\]
\end{ejercicio}

\begin{ejercicio}
Sea $f:[0,1]\rightarrow \R$ una función tal que, para todo $x\in[0,1]$, 
$\int_{0}^{x}f dx = \int_{x}^{1}f dx$. Demuestre que $f(x)=0$ para todo
$x \in [0,1]$.
\end{ejercicio}

\begin{ejercicio}
Sean $f,g$ funciones integrables $[a,b]$. Demuestre que la función 
$M(x)=max(f(x),g(x))$ es integrable en $[a,b]$
\end{ejercicio}

%% ============================================================
\part{Análisis Funcional}

\chapter{Primer Parcial}

\section{Examen}

\begin{ejercicio}
Si $(X, d)$ es un espacio metrico, demuestre que $\bar{d}(x,y)=\frac{d(x,y)}{1+d(x,y)}$
es otra métrica en $X$.
\end{ejercicio}

\begin{proof}
  i)\\
  Como $d$ es métrica en $X$ $\Rightarrow d(x,y) \geq 0 \quad \forall x,y \in X$\\
  $\therefore \bar{d}(x,y)\geq \quad \forall x,y \in X$\\
  ii)\\
  Como $d(x,y)=0 \iff x=y$\\
  $\therefore \bar{d}(x,y)=\frac{d(x,y)}{1+d(x,y)}=0 \iff x=y$\\
  iii)\\
  Como $d(x,y)=d(y,x)\quad \forall x,y \in X$\\
  $\Rightarrow 1+ d(x,y)=1+d(y,x)$\\
  $\therefore \frac{d(x,y)}{1+d(x,y)}=\frac{d(y,x)}{1+d(y,x)} \quad \forall x,y \in X$\\
  iv)\\
  Sea $f(x)=\frac{x}{1+x}$ con $x\in\R$, $f'(x)=\frac{1}{(1+x)^2}$\\
  entonces $f$ es diferenciable y $f'(x)>0 \quad \forall x\in \R$ que por teorema,
  $f$ es estrictamente creciente\\
  $d$ es métrica así que $d(x,z)\leq d(x,y)+d(y,z)\quad x,y,z \in X$ y al ser $f$ estrictamente creciente,
  $f(d(x,z))\leq f(d(x,y)+d(y,z))$ $\Rightarrow \bar{d}(x,z)=\frac{d(x,z)}{1+d(x,z)}\leq \frac{d(x,y)+d(y,z)}{1+d(x,y)+d(y,z)}$\\
  $\bar{d}(x,z)=\frac{d(x,y)}{1+d(x,y)+d(y,z)}+\frac{d(y,z)}{1+d(x,y)+d(y,z)}$\\
  Como $d(x,y)\geq 0$ y $d(y,z)\geq0$; $1+d(x,y)+d(y,z)\geq 1 + d(x,y)$ y $1+d(x,y)+d(y,z)\geq 1 + d(y,z)$\\
  Luego multiplicamos $d(x,y$) por $\frac{1}{1+d(x,y)+d(y,z)}\leq \frac{1}{1+d(x,y)}$; 
  $\frac{d(x,y)}{1+d(x,y)+d(y,z)}\leq \frac{d(x,y)}{1+d(x,y)}$\\
  De forma análoga, $\frac{d(y,z)}{1+d(x,y)+d(y,z)}\leq \frac{d(y,z)}{1+d(y,z)}$\\
  De manera que ahora podemos decir $\bar{d}(x,z)\leq \frac{d(y,z)}{1+d(y,z)} + \frac{d(x,y)}{1+d(x,y)}=\bar{d}(x,y)+\bar{d}(y,z)$ 
  $\forall x,y,z \in X$\\
  $\therefore$ se cumplen los axiomas para determinar que $\bar{d}$ es métrica.
\end{proof}

\begin{ejercicio}
  Un conjunto $A$ en un espacio métrico $(X,d)$ es acotado si 
  $\delta(A)=sup\conj{d(x,y)|x,y\in A}<\infty$. Demuestre que la unión de dos 
  conjuntos acotados es acotado.
\end{ejercicio}

\begin{proof}
  Sean $A,B$ conjuntos acotados.\\
  Si $A=\emptyset$ o $B=\emptyset$, $A\cup B$ es acotado. Suponga entonces $A,B \not = \emptyset$\\
  Suponga $A\cup B$ no es acotada ergo, $sup\conj{d(x,y)|x,y\in A\cup B}=\infty$\\
  Entonces para cada $n\in \N$ existen $x_n, z_n \in A\cup B$ con $d(x_n,z_n) \geq n$\\
  Como $A$ y $B$ son acotados, tomamos $N\in \N$ con $N>\delta(A)$ y $N>\delta(B)$\\
  Si $n\geq N$, $x_n$ y $z_n$ no puden ser del mismo conjunto:\\
  $x_n, z_n ,\in A \Rightarrow \quad d(x_n,z_n)\geq n>N> \delta(A)=sup\conj{d(u,v)|u,v\in A}$\\
  $x_n, z_n ,\in B \Rightarrow \quad d(x_n,z_n)\geq n>N> \delta(B)=sup\conj{d(u,v)|u,v\in B}$\\
  En ambos casos se llegan a una contradicción por la definición del supremo\\
  Suponga $x_n \in A, z_n\in B \forall n \geq N$, pues al tener 
  $d(x,y)=d(y,x) \forall x,y \in X$ podemos intercambiar $x_n$ con $z_n$ si $x_n \in B$ y $z_n \in A$
  para algunos $n \geq N$\\
  Sea $y\in A$ $d(x_n,z_n)\leq d(x_n,y)+ d(y,z_n)$ Tomando $n\rightarrow \infty$ como 
  $d(x_n,z_n)\rightarrow \infty \Rightarrow \quad d(x_n,y)+ d(y,z_n)\rightarrow \infty$\\
  Pero $x_n,y \in A \Rightarrow \quad d(x_n, y)\leq \delta(A) \quad \forall n\geq N$, por lo que 
  $d(y,z_n)\rightarrow \infty$ cuando $n\rightarrow \infty$\\
  Además, si $w\in B$ fijo, $d(y,z_n)\leq d(y,w)+d(w,z_n)$\\
  $\Rightarrow d(y,z_n)\leq d(y,w)+\delta(B) \quad \forall n\geq N$\\
  Contradiciendo que $d(y,z_n)$ no está acotado. Por tanto, $A\cup B$ es acotado.
\end{proof}

\begin{ejercicio}
  Sea $x_0$ un punto fijo en el espacio métrico $(X,d)$ y $r>0$. Demuestre que 
  $\conj{x\in X | d(x,x_0)\leq r}$ es cerrado. 
\end{ejercicio}

\begin{proof}
  Llamemos $\bar{B_r}(x_0)$ al conjunto dado y probemos que es cerrado mostrando que 
  $X-\bar{B_r}(x_0)$ es abierto.\\
  Tomemos $y_0 \in X- \bar{B_r}(x_0)$\\
  Como $y_0 \not \in \bar{B_r}(x_0)\Rightarrow \quad d(x_0,y_0)>r$\\
  Consideremos $r'=d(x_0,y_0) - r>0$ y tomemos $y\in \bar{B_{r'}}(y_0)$\\
  Notemos que $d(y,y_0)<d(x_0,y_0)-r\leq d(y_0,y)+d(y,x_0)-r$\\
  $\Rightarrow r < d(y,x_0)$\\
  Así que $y\not \in \bar{B_r}(x_0)$\\
  $\Rightarrow X-\bar{B_r}(x_0)$ es abierto\\
  $\therefore \bar{B_r}(x_0)$ es cerrado
\end{proof}

\begin{ejercicio}
  Si $(x_n)$ y $(y_n)$ son sucesiones de Cauchy en un espacio métrico $(X,d)$, 
  demuestre que $(d(x_n,y_n))$ converge
\end{ejercicio}

\begin{proof}
Como $(x_n)$ es Cauchy \\
$\Rightarrow \forall \epsilon>0 \quad \exists N_1\in \N |$ si $n,m>N_1\Rightarrow d(x_n,x_m)<\frac{\epsilon}{2}$\\
Y como $(y_n)$ es Cauchy \\
$\Rightarrow \forall \epsilon>0 \quad \exists N_2\in \N |$ si $n,m>N_2\Rightarrow d(y_n,y_m)<\frac{\epsilon}{2}$\\
Sea $N_0=max\conj{N_1,N_2}$\\
$\Rightarrow d(x_n,y_n)\leq d(x_n, x_m) + d(y_n, y_m)+ d(x_m,y_m)$\\
$\Rightarrow d(x_n,y_n)-d(x_m,y_m)\leq d(x_n, x_m) + d(y_n, y_m)$\\
$\Rightarrow d(x_n,y_n)-d(x_m,y_m) < \frac{\epsilon}{2}+\frac{\epsilon}{2}=\epsilon$\\
Estas desigualdades se pueden realizar de manera análoga para $d(x_m,y_m)-d(x_n,y_n) < \epsilon $\\
$\Rightarrow |d(x_n,y_n)-d(x_m,y_m)| < \epsilon $ $\forall \epsilon>0 $ y $n,m>N_0$\\
$\therefore \conj{d(x_n,y_n)}$ es una sucesión de Cuachy en $\R$ con la métrica usual, la cual sabemos es un espacio
completo\\
$\therefore \conj{d(x_n,y_n)}$ es una sucesión convergente
\end{proof}

\begin{ejercicio}
  Sea $X =\R^n$ con métrica $d(x,y)=max\conj{|x_i-y_i| \quad | i=1,2,...,n}$
  donde $x=(x_1,x_2,...,x_n)$, $y=(y_1,y_2,...,y_n)$. Demuestre que $(X,d)$
  es completo.
\end{ejercicio}

\begin{proof}
  Tomeos $(x_n)$ sucesión de Cauchy en $\R^n$\\
  $\Rightarrow \forall \epsilon>0 \quad \exists N_0\in \N |$ si $n,m>N_0\Rightarrow d(x^n,x^m)<\epsilon$\\
  $\Rightarrow max\conj{|x_i^n-x_i^m| \quad | i\in \conj{1,...,n}}<\epsilon$\\
  O sea que la máxima diferencia entre todas las componentes de dos elementos de la sucesión es tan chica como queramos
  y si consideramos $i\in \conj{1,...,n}$\\
  $\Rightarrow \forall \epsilon>0 \quad \exists N_i\in \N |$ si $n,m>N_i\Rightarrow d(x_i^n,x_i^m)<\epsilon$\\
  Esto nos da una sucesión de Cauchy en $\R$, el cual es completo\\
  $\therefore (x_i^{n})_{n \in \N}$ es sucesión convergente\\
  Esto $\forall i\in \conj{1,...,n}$ converge a un elemento de $\R$, digamos $k_i \Rightarrow k=(k_1,k_2,...,k_n)\in \R^n$\\
  Sea $N=max\conj{N_i | i\in \conj{1,...,n}}$\\
  $\forall \epsilon >0$ si $n>N \Rightarrow \quad d(x^n,k)<\epsilon$\\
  $\therefore (x^{n})$ converge a un elemento de $X$, y como Cauchy es una arbitraria\\ 
  $\therefore (X,d)$ es un espacio completo
\end{proof}

\chapter{Segundo Parcial}

\section{Examen}

\begin{ejercicio}
  Si $(X_1, \|.\|_1)$ y $(X_2, \|.\|_2)$ son espacios normados, demuestre que el 
  espacio vectorial $X=X_1\times X_2$ es un espacio normado si se define 
  $\|x\|=max\conj{\|x_1\|_1,\|x_2\|_2}$ donde $x=(x_1,x_2)$.
\end{ejercicio}

\begin{ejercicio}
  Si dos normas $\|.\|$ y $\|.\|_0$ en un espacio vectorial $X$ son equivalentes, demuestre
  que $\|x_n-x\|\rightarrow 0 $ implica $\|x_n-x\|_0\rightarrow 0$.
\end{ejercicio}

\begin{ejercicio}
  Un espacio métrico $X$ es localmente compacto si cada punto de $X$ tiene una vecindad 
  compacta. Demuestre que un espacio métrico compacto $X$ es localmente compacto.
\end{ejercicio}

\begin{ejercicio}
  Si $X$ es un espacio métrico compacto y $M \subseteq X$ es cerrado, demuestre que $M$ 
  es compacto.
\end{ejercicio}

\begin{ejercicio}
  Si $X$ y $Y$ son espacios vectoriales de dimensión finita y $T:X\rightarrow Y$ (lineal) es uno a uno.
  Demuestre que $dimX\leq dimY$
\end{ejercicio}

\chapter{Tercer Parcial}

\section{Examen}

\begin{ejercicio}
  Si $x$ y $y$ son vectores diferetes en un espacio vectorial $X$ de dimensión finita. Demuestre 
  que existe un funcional lineal $f$ en $X$ tal que $f(x) \not = f(y)$.
\end{ejercicio}

\begin{ejercicio}
  Si $X$ es un espacio normado y $dim X = \infty$, demuestre que el espacio dual 
  $X'$ no es idéntico al espacio dual algebraico $X*$.
\end{ejercicio}

\begin{ejercicio}
  Demuestre que dos funcionales lineales $f_1 \not = 0$ y $f_2 \not = 0$
  definidos en el mismo espacio vectorial tales que sus espacios nulos son iguales,
  son proporcionales.
\end{ejercicio}

\begin{ejercicio}
  Sea $f_1 \not = 0$ un funcional lineal en el espacio vectorial $X$ y $x_0$ un elemento fijo 
  de $X-N(f)$. Demuestre que $\forall x \in X$, $x=\alpha x_0 + y$ donde $y \in N(f)$.
\end{ejercicio}

\begin{ejercicio}
  Sea $T:l^{\infty}\rightarrow l^{\infty}$ definido por $T_x=(\frac{x_i}{i})$ donde $x=(x_i)$.
  Demuestre que $R(T)$ no es cerrado.
\end{ejercicio}

%% ============================================================
%%  ÍNDICE ANALÍTICO
\printindex

\end{document}
